\subsection{The selected-cusp source loop after its first fixed-edge flop}
\begin{statement}{P3K.40}{Refuted; physical-line correction}\par\noindent\textbf{Physical-line correction, 17 November 2026.} The actual first fixed-edge return is \(-7=21\), replacing the old \(-5=23\). Its native right integer \(-6\) and every local source/Cox/cap comparison remain valid. The initial physical class is \(-1\), and the same directed flop adds \(-6\). K68--69 prove the correction. The following earlier assertion and its argument are retained as the located historical claim; their corrected current consequences are proved in K68--69.\medskip\par 
Retain the actual positive selected-cusp source loop of P3F.54 and
P3F.56, with fixed sufficiently small scalar, the III-star cap, no
logarithmic modification at either member of the coalescing finite
pair, and the initial physical-delta-edge fan. Its parameter is one
positive \(\rho\)-loop with \(a(u)=\rho^2\), the actual twofold
marked closure of the selected source puncture. Flop that fixed edge
in the periodic cusp completion, carrying its actual source, line,
cap and ordinary-annulus maps. The positive return of the resulting
family has oriented homotopy-sphere class
\[
                 [\Sigma_{\Phi,\mathrm{flop}}]
                     =-5=23\quad\hbox{in }\mathbb Z/28.
\]
In the actual first-resolution relative product, its initial right
exponent is \(n=-6\). The other free exponent \(m\) is retained
without assigning its value. This statement assigns neither the
separate Jin characteristic pair nor a return for another edge or
an arbitrary fan.
\end{statement}
\begin{proof}
\par

\medskip\noindent\textbf{SF.1. The closed product is supplied by the actual source coordinates}

\par\nobreak

Write the trace of the fixed curve, over a small punctured rho-disc, as
\(F=\mathbf P^1_s\times D_\rho^*\). Its two distinguished sections are
\(D_+=\{s=\rho\}\) and \(D_-=\{s=-\rho\}\). The literal ordered
coordinate and inverse are
\[
 Z=\frac{s-\rho}{s+\rho},\qquad
 s=\rho\frac{1+Z}{1-Z}.                                      \tag{SF.1}
\]
This is an actual biholomorphism from \(F\) to
\(\mathbf P^1_Z\times D_\rho^*\). Its determinant is \(2\rho\), so
it has no zero on this base. In particular it is single-valued after
one rho turn, and retains the two ordered endpoint sections.

The relative small-resolution charts used in P3F.56 have normal
coordinates \(c_0,c_1\), with the actual overlap
\[
             c_i'=Zc_i,\qquad i=0,1,                       \tag{SF.2}
\]
and physical delta weights \((-1,+1)\). Thus their normal bundle is the
two copies of \(\mathcal O_{\mathbf P^1}(-1)\) in precisely the
coordinate (SF.1). This is also the Cartier identification used in
P3F.56: after conjugating the negative-weight normal coordinate its
spinor is
\(\mathcal O(D_+)\oplus\mathcal O(-D_-)\), and the determinant
trivialization is the ratio \(Z\). We use that same identification.

Here is why no source-scaling loop is discarded in using (SF.2).
The source chart and its actual Tate parameter satisfy
\[
 \xi\eta=q,\qquad
 \frac{q_T}{q}=-\frac{u(\rho)}{64\rho^6}+O(q).               \tag{SF.3}
\]
The coefficient in (SF.3) has winding \(-6\) on the rho-circle.
Choose the source collar uniformly smaller than its distance from the
second moving critical value; then this coefficient is an actual
nonzero holomorphic function on the whole parameter annulus and
normal coordinate polydisc. Every remaining nonzero toric chart unit
\(U\) on such a representative has an exact factorization
\[
                    U=\rho^{e}\exp H,\qquad e\in\mathbb Z. \tag{SF.4}
\]
Indeed \(e\) is its integral winding on the rho generator; subtracting
\(e\,d\rho/\rho\) from \(dU/U\) kills its only possible period,
and integration gives the single-valued holomorphic \(H\). This
argument is on the annulus, not across its missing centre. The integer
\(-6\) in (SF.3) is retained as one of these exponents.

The source and line units in P3F.56 extend on the original source
polydisc. Their changes to the Tate cell use (SF.3), so the preceding
factorization is applied only after that exact coordinate change.
The period action has trivial stabilizer on an individual conifold
cell. We choose one representative of the cell containing the fixed
edge and transport its charts to all translates. Thus (SF.4) and its
chosen logarithm give the same maps on both resolution charts and
on their actual ordinary annulus. Contracting \(H\) and the positive
modulus transports their common truncation collar by the resulting
isotopy; it leaves the integral angular powers in place. The source
coordinate (SF.1) is carried along, rather than identified with the
old source coordinate by an unmarked homotopy.

Equivalently, for the linear normal maps this assertion can be checked
over the whole compact fixed curve. Delta-equivariance preserves its
two distinct weight lines. After (SF.1)--(SF.2), an automorphism of
either line is multiplication by a holomorphic function on
\(\mathbf P^1\), hence by a function of rho alone. Its integer winding
is still retained by (SF.4). Higher normal terms are removed by the
literal equivariant normal rescaling and isotopy extension on the
same collar. The linearization homotopy has invertible diagonal
blocks throughout and starts with the actual map. No change of an
outer gluing or of an initial normal frame is made.

\par

\medskip\noindent\textbf{SF.2. The two relative family products have the same boundary arrow}

\par\nobreak

The common conifold link is
\[
 (S^3_a\times S^3_c)/S^1,\qquad
 (a,c)\sim(e^{i\phi}a,e^{-i\phi}c).
\]
Put, in the quaternion convention of P3K.23,
\[
 u=a_0+ja_1,\qquad v=c_0+j\overline{c_1},\qquad
 q_H=uv,\quad h=uiu^{-1},\quad k=v^{-1}iv.                 \tag{SF.5}
\]
The map to the first product is \([a,c]\mapsto(h,q_H)\), and that to
the second product is \([a,c]\mapsto(k,q_H)\). Their inverses choose
one Hopf representative and use respectively \(v=u^{-1}q_H\) or
\(u=q_Hv^{-1}\); changing the representative is exactly the displayed
auxiliary circle. These are the actual normal-sphere products of the
two small resolutions.

On the first source chart (SF.5) is precisely P3F.56's retained map
\[
              q_H=(1+jZ)(c_0+j\overline{c_1}),             \tag{SF.6}
\]
after the common positive normalization of its norm. The positive
normalizing factor extends across both charts and is carried along
in both resolutions. Thus the local identification in this proof is
the one in the proved trace comparison, including its normal frame.

The exact common boundary map, computed from the same two factors,
is
\[
               \tau(h,q_H)=(q_H^{-1}hq_H,q_H).             \tag{SF.7}
\]
It has no rho dependence. To verify this for all retained integer
chart powers, give the four Cox coordinates \((a_0,a_1,c_0,c_1)\)
their actual integer weights \((A_0,A_1,C_0,C_1)\). Such integer
lifts exist because the auxiliary-circle vector \((1,1,-1,-1)\)
is primitive; the quotient cocharacter lattice is free. We impose
no condition on their sum. Put
\[
 \begin{aligned}
 a_\theta&=e^{i(A_0-A_1)\theta/2},&
 d_\theta&=e^{i(A_0+A_1)\theta/2},\\
 e_\theta&=e^{i(C_0+C_1)\theta/2},&
 b_\theta&=e^{i(C_0-C_1)\theta/2},\\
 S&=A_0+A_1+C_0+C_1.
 \end{aligned}
\]
The literal quaternion changes are
\(u'=a_\theta u d_\theta\) and
\(v'=e_\theta v b_\theta\). Therefore
\[
 \begin{split}
 h'&=a_\theta h a_\theta^{-1},\qquad
 k'=b_\theta^{-1}k b_\theta,\\
 q_H'&=a_\theta u e^{iS\theta/2}v b_\theta\\
     &=a_\theta e^{hS\theta/2}q_H b_\theta
       =a_\theta q_H e^{kS\theta/2}b_\theta .
 \end{split}                                             \tag{SF.8}
\]
Here \(h=u i u^{-1}\), \(k=v^{-1}iv\), and both displayed
exponential identities follow by conjugating the same middle
complex scalar. Substitution into (SF.7) gives
\[
 (q_H')^{-1}h'q_H'
 =b_\theta^{-1}e^{-kS\theta/2}
              (q_H^{-1}h q_H)e^{kS\theta/2}b_\theta
 =b_\theta^{-1}k b_\theta=k'.
\]
This proves equivariance even when \(S\ne0\). At \(\theta=2\pi\)
the factors \(a_\theta,d_\theta\) are central signs whose product
is \((-1)^{2A_0}=1\), and the product of the two signs
\(e_\theta,b_\theta\) is \((-1)^{2C_0}=1\). Thus both actual Cox
factors, and consequently both product actions, close. The possible
nonzero middle exponential has been retained.

Formula (SF.8) consequently retains all integer windings from
(SF.3)--(SF.4), including the actual sixth source scaling. It does not
require evaluating or discarding any of them. The same comparison
holds during the common unit isotopy in SF.1, so its initial and
final ordinary-annulus maps are the original ones. The earlier
working passage with \(S=0\) is preserved as superseded in
\path{code/part3/attempts/source_cusp_fixed_edge_sum_zero_restriction/};
that restricted passage was not promoted as a theorem.

Choose the closed product connection just described on the fixed
curve tube and equivariant connections on its complement; an
invariant partition of unity patches them while retaining the fixed
collar. The resulting physical return is the identity near that
curve. The same choice in the flopped family gives the identical
common-exterior return, conjugated only by the extension of (SF.7),
\[
 T(y,q_H)=(q_H^{-1}yq_H,q_H),\qquad
 T^{-1}(y,q_H)=(q_Hyq_H^{-1},q_H).                         \tag{SF.9}
\]
It is smooth at \(y=0\). The first and second normal conventions,
their complex orientations and their physical Hopf circle are exactly
those in P3K.23. This verifies the family/framing hypothesis required
for that statement, for the source loop itself.

\par

\medskip\noindent\textbf{SF.3. The actual relative integer and the flopped class}

\par\nobreak

Write the actual first relative return in the proved relative group
as \(L^mR^n\mathcal C^s\), using P3K.23's maps
\[
 L(y,q_H)=(y,\zeta(y)q_H),\qquad
 R(y,q_H)=(y,q_H\zeta(y)),
\]
and its positively oriented central generator \(\mathcal C\).
We retain the actual, presently unevaluated integer \(m\).
The central coordinate is the oriented closed six-sphere return,
so P3F.56 gives \(s=1\pmod {28}\).

The coefficient \(n\) is fixed without a separate normal
Pontryagin calculation. We spell out the comparison being used.
The first quaternion normal product has physical circle
\(q_H\mapsto q_He^{-i\theta}\), so the exact Hopf quotient has
orientation \(S^2_z\times S^3_y\). Its relative variation is
\(-k_{\rm Jin}\). The actual generator \(L\) has quotient
characteristic pair \((-1,-4)\), computed from its adjoint clutch
with \(p_1=-4\). Removing its variation therefore changes the
characteristic pair to \((0,\ell-4k)\). The remaining free relative
map is \((LR^{-1})^b\). On its fixed three-sphere its normal clutch
is \(\operatorname{Ad}_{\zeta^b}\), hence its Pontryagin evaluation
is \(-4b\). The integral Gysin pullback is injective on the relevant
fourth cohomology and the actual marked cycle evaluates to one.
Consequently, in these same products,
\[
           \ell-4k=-4b,
           b=6\frac{4k-\ell}{24},\qquad
           n=-b=-6A.                                    \tag{SF.10}
\]
Every arrow and orientation in this argument is the one verified in
P3K.23, now applicable by SF.1--SF.2; it does not import the scalar-loop
value of either \(k\) or \(\ell\). Since the actual source loop has
\(A=1\) by P3F.56, it gives \(n=-6\).

The literal conjugation in (SF.9) is
\[
 T(L^mR^n\mathcal C^s)T^{-1}
       =L^{-n}R^{m+2n}\mathcal C^{s+n}.                   \tag{SF.11}
\]
It follows that the actual positive source-loop return after its
first fixed-edge flop has class
\[
                          1-6=-5=23\pmod {28}.           \tag{SF.12}
\]
This computation retains the actual source parameter, fixed scalar,
cap norm, physical delta character, and ordinary annular maps. It
does not replace the source loop by a scalar loop.

The transported complex first-resolution reset of P3K.31 changes the
relative pair to \((m-12,6)\), while retaining the same central class
\(-5\). To compute a later flop at another edge one still needs the
source loop's full common-exterior variation and derivative. Neither
\(m\), the separate Jin pair, nor an arbitrary-fan formula is assigned
by the present proof.
\end{proof}
\par

\medskip\noindent\textbf{Audit and computation status}

\par\nobreak

Root has independently passed the full SF.1--3 proof and the exact
K23 Hopf/relative-group comparison. The independent SF.1--2 audit also
passes the actual source-coordinate, normal-line and unrestricted
four-weight quaternion maps. The exact
source-coordinate, quaternion equivariance and relative-group checks
are executed by \path{code/part3/source_cusp_fixed_edge.py}, with the matching
JSON and captured output. No result from an earlier run is used.
